\section{Related Work}
\label{sec:related}


\paragraph{Residual Stream Magnitude Scaling.}
Residual connections preserve input representations while allowing each block
to add a learned transformation~\citep{he2016deep}. In transformers, this creates
a residual stream that carries information across layers and serves as the shared
workspace for attention and feed-forward updates~\citep{elhage2021mathematical}.
This preservation path is also a primary target for architectural scaling:
mechanisms such as ReZero~\citep{Bachlechner2020ReZeroIA}, LayerScale~\citep{touvron2021going},
and deep-transformer normalization schemes~\citep{xiong2020on,Wang2022DeepNetST}
adjust branch scale to stabilize optimization, while residual paths mitigate
representational degeneracy such as rank collapse~\citep{dong2021attention}.
These works establish the importance of residual pathways and branch magnitude
for training stability. In contrast, we study a complementary directional geometric property of the learned update itself: whether it reinforces the incoming representation
direction or steers it into orthogonal semantic subspaces. 

\paragraph{Self-Directed Attention Value Flow.}
Attention explicitly aggregates information across
tokens~\citep{vaswani2017attention}, but attention mass can also concentrate
on special or persistent positions, as shown by attention-sink behavior in
long-context inference~\citep{xiao2023efficient}. Gated attention variants
mitigate such sink behavior by modulating attention
outputs~\citep{qiu2025gatedattention}. However, even when attention sinks are reduced,
large diagonal attention weights can still route substantial mass back to the
current token itself. Exclusive Self-Attention (\xsa{}) directly addresses this
self-directed pathway by modifying value flow aligned with the current
token~\citep{zhai2025xsa}. In this work, we analyze this self-directed pathway from a directional geometric
perspective, disentangling direction-preserving updates from direction-changing
information. 
